\usepackage{keytheorems}
\usepackage[nameinlink]{cleveref}

%% Helpers for theorem styles and environments.
%% NOTE: cleveref doesn't work well with \newtheorem{...}[theorem]{...},
%%   so we use aliascnt to create theorem environments that share the same counter as the main theorem environment.
\newcounter{thmenvcnt}[chapter]
\renewcommand{\thethmenvcnt}{\thechapter.\arabic{thmenvcnt}}

% Create a new theorem environment that shares the counter.
%   #1 - the name of the new environment
%   #2 - the title of the new environment
%   #3 - (optional) the ending symbol for the environment (defaults to empty)
%   #4 - (optional) the cref name (defaults to \MakeLowercase{#2})
%   #5 - (optional) the cref plural name (defaults to #4 + 's')
%   #6 - (optional) the Cref name (defaults to #2)
%   #7 - (optional) the Cref plural name (defaults to #6 + 's')
\NewDocumentCommand{\newtheoremalias}{m m O{} O{\MakeLowercase{#2}} O{#4s} O{#2} O{#6s}}{%
  \newkeytheorem{#1}[
    title={#2},
    sibling=thmenvcnt,
    refname={#4,#5},
    Refname={#6,#7},
    qed={#3} 
  ]
  \newkeytheorem{#1*}[
    title={#2},
    numbered=no,
  ]
}

%% Theorem styles.
\newtheoremstyle{sans-plain}
  {\topsep}
  {\topsep}
  {\itshape}
  {}
  {\sffamily\bfseries}
  {.}
  {.5em}
  {}
\newtheoremstyle{sans-definition}
  {\topsep}
  {\topsep}
  {\normalfont}
  {}
  {\sffamily\bfseries}
  {.}
  {.5em}
  {}

%% Theorem environments.
\theoremstyle{sans-plain}
\newtheoremalias{theorem}{Theorem}
\newtheoremalias{lemma}{Lemma}
\newtheoremalias{proposition}{Proposition}
\newtheoremalias{corollary}{Corollary}[][corollary][corollaries][Corollary][Corollaries]

\theoremstyle{sans-definition}
\newtheoremalias{definition}{Definition}
\newtheoremalias{example}{Example}
\newtheoremalias{remark}{Remark}[$\lrcorner$]
\newtheoremalias{question}{Question}
\newtheoremalias{acknowledgement}{Acknowledgement}

% todo notes
\newcommand{\todo}[1]{%
  \PackageWarning{TODO}{#1}%
  \vskip0.5em%
  \noindent\textcolor{red}{\textbf{TODO.}~#1}%
  \vskip0.5em%
}

%% Cross-reference styles.
\NewCommandCopy{\crefOrig}{\cref}
\NewCommandCopy{\CrefOrig}{\Cref}
\renewcommand{\cref}[1]{\textsf{\crefOrig{#1}}}
\renewcommand{\Cref}[1]{\textsf{\CrefOrig{#1}}}
